\subsection{Spectrum of an endomorphism}

DEFINITION 1. --- Let E be a topological vector space and let u be an endomorphism of E. We call the spectrum of u, and denote by Sp(u), the spectrum of u relative to the unital algebra \(\mathcal{L}(\mathrm{E})\) .

Let E be a topological vector space and let \(u \in \mathcal{L}(\mathrm{E})\) . The spectrum of u is the set of complex numbers \(\lambda\) such that \(u - \lambda1_{E}\) is not an automorphism of E. If E is metrizable and complete, it is also the set of complex numbers \(\lambda\) such that \(u - \lambda1_{E}\) is not bijective (EVT, I, p.~19, cor. 1).

Every eigenvalue of u belongs to the spectrum of u, but the converse is false in general.

In the rest of this paragraph, we shall restrict ourselves to studying the notion of spectrum in the case where E is a Banach space.

Lemma 1. --- Let E be a Banach space and let u be an endomorphism of E. Let \((\mathrm{E}_{i})_{i\in\mathrm{I}}\) be a finite family of closed subspaces of E, stable under u, such that \(E = \bigoplus_{i\in I} E_{i}\) . For every \(i \in I\) , let \(u_{i}\) be the endomorphism of \(E_{i}\) induced by u. We have \(\mathrm{Sp}(u) = \bigcup_{i\in I} \mathrm{Sp}(u_{i})\) , and for every \(f \in \mathcal{O}(\mathrm{Sp}(u))\) , the endomorphism \(f(u)\) stabilizes the spaces \(E_{i}\) , and \(f(u)\) coincides with \(f(u_{i})\) on \(E_{i}\) .

The endomorphism \(u\) is an isomorphism if and only if \(u_i\) is an isomorphism for every \(i \in \mathrm{I}\) . Applying this property to \(u - \lambda 1_{\mathrm{E}}\) , one deduces that \(\operatorname{Sp}(u)\) is the union of the sets \(\operatorname{Sp}(u_i)\) for \(i \in \mathrm{I}\) .

Let \(f \in \mathcal{O}(\mathrm{Sp}(u))\) . The endomorphism \(f(u)\) of E belongs to the bicommutant of \(u\) in \(\mathcal{L}(\mathrm{E})\) (theorem 5 of I, p.~74), hence it commutes with the projections \(p_i\) . It then stabilizes the spaces \(\mathrm{E}_i\) . Consider the continuous unital morphism \(\varpi\) of the product algebra \(\prod_{i \in \mathrm{I}} \mathcal{L}(\mathrm{E}_i)\) into \(\mathcal{L}(\mathrm{E})\) defined by \((v_i)_{i \in \mathrm{I}} \mapsto \bigoplus_i v_i\) . It maps the family \((u_i)\) to \(u\) . Proposition 7 of I, p.~75 then implies that \(\mathrm{Sp}(u) \subset \bigcup_{i \in \mathrm{I}} \mathrm{Sp}(u_i)\) and

\[
f (u) = f (\varpi ((u _ {i}) _ {i \in \mathrm{I}})) = \varpi ((f (u _ {i})) _ {i \in \mathrm{I}}) = \bigoplus_ {i \in \mathrm{I}} f (u _ {i}),
\]

which concludes the proof of the lemma.

PROPOSITION 1. --- Let E be a complex Banach space and let u be an endomorphism of E. Let \(\lambda \in \mathbf{C}\) and \(f \in \mathcal{O}(\mathrm{Sp}(u))\) . We have

\[
\operatorname{Ker} (u - \lambda 1 _ {\mathrm{E}}) \subset \operatorname{Ker} (f (u) - f (\lambda) 1 _ {\mathrm{E}}).
\]

Let \(x \in E\) be nonzero. The set A of \(v \in \mathcal{L}(E)\) such that \(x\) is an eigenvector of \(v\) is a unital subalgebra of \(\mathcal{L}(E)\) .

The algebra A is full: if \(v \in A\) is invertible in \(\mathcal{L}(\mathrm{E})\) and if x is an eigenvector of v for the eigenvalue \(\lambda\) , then \(\lambda \neq 0\) and \(v^{-1}(x) = \lambda^{-1}x\) which proves that \(v^{-1} \in A\) .

The algebra A is closed in \(\mathcal{L}(\mathrm{E})\) . Indeed, if \((v_n)_{n\in \mathbf{N}}\) is a sequence in A such that \(v_{n}\) converges to \(v\in \mathcal{L}(\mathrm{E})\) , the sequence \((\lambda_{n})_{n\in \mathbf{N}}\) such that \(v_{n}(x) = \lambda_{n}x\) is bounded, hence admits a subsequence converging to a complex number \(\mu\) , so that \(v(x) = \mu x\) .

Suppose that \(x\) is an eigenvector of \(u\) and that \(u(x) = \lambda x\) . The algebra A contains \(u\) , hence contains the closed unital full subalgebra B generated by \(u\) , which is commutative. The map \(\chi \colon \mathrm{B} \to \mathbf{C}\) which associates with \(v\) the eigenvalue of \(v\) relative to \(x\) is a character of B such that \(\chi(u) = \lambda\) . For every \(f \in \mathcal{O}(\mathrm{Sp}(u))\) , one has \(f(u) \in \mathrm{B}\) and \(\chi(f(u)) = f(\chi(u)) = f(\lambda)\) by prop. 7 of I, p.~75, whence the proposition.

